% ============================================================================
\clearpage
\section*{v4.9.0-r06 Source-Only: The Tachyon Hit Chain, its Acceleration, and the Higgs Field}
\addcontentsline{toc}{section}{v4.9.0-r06 Source-Only: The Tachyon Hit Chain, its Acceleration, and the Higgs Field}
\label{sec:v490-r06-hit-chain-higgs}
% ============================================================================

\paragraph{Purpose.}
Theoremlet r05.3(c) stated that a photon returning to the origin of a tachyon cannot create it.  The author objected that the
tachyon already runs backward in time, and then refined the picture.  The Higgs field tries to create matter along a line.
A tachyon that has been travelling along this line for a long time hits it again and again.  Each hit is a point of
spacetime at which an essentially identical photon arises.  This layer makes r05.3(c) precise, works out the hit chain,
asks what an accelerating or infinitely accelerated tachyon does to the hits, and checks the picture against the Higgs
field.  Special relativity is used with \(c=1\).  The computations are \texttt{llm/k42\_tachyon\_trefferkette.py}.  No
value is compared with data.  Earlier layers are not rewritten.

\noindent\textbf{Source-only status.}
This v4.9.0-r06 layer is source-only.  It does not build a PDF, does not assign a new DOI,
and does not authorize a public release.  The v4.9.0 Version DOI
\href{https://doi.org/10.5281/zenodo.23534954}{10.5281/zenodo.23534954} and the Concept DOI
\href{https://doi.org/10.5281/zenodo.19210728}{10.5281/zenodo.19210728} are unchanged by it.

\subsection*{v4.9.0-r06 closing the loop, made precise}
\label{subsec:v490-r06-loop}

\paragraph{Theoremlet v4.9.0-r06.1.}
\begin{enumerate}[leftmargin=2em]
\item[(a)] A photon needs the time \(d\) for the distance \(d\).  A tachyon of speed \(v>1\) covers the same distance in the
time \(d/v\), and even read backward it takes back only \(d/v<d\).  One photon leg and one tachyon leg therefore never close
a loop.  The remaining gap is \(d(1-1/v)\), which is \(0.451\,d\) for \(v=1/\bsym\).
\item[(b)] Two tachyon legs emitted in frames moving relative to each other can close or overshoot a loop.  A return leg
that is instantaneous in a frame of speed \(u\) arrives at the origin at \(d/v-ud\), before the start for \(u>1/v\).
\item[(c)] For \(v=1/\bsym\) and \(u=\bsym\) the return arrives exactly at the start.  It runs back along the same
spacelike line, because the outgoing leg is itself instantaneous in that frame.  This is one exchange read in two
directions, case (a) of Theoremlet r02.1, and carries no information around a loop.
\end{enumerate}

\paragraph{Proof sketch.}
Arithmetic on the worldlines; the return leg is the line of simultaneity \(t-t_B=u(x-x_B)\) through
\(B=(d/v,d)\).  Checked numerically.  \(\square\)

\paragraph{Correction of wording in r05.}
The phrase ``a photon leg running backward in time'' in the reading of r05.3 is made precise by (a) and (b).  What fails
is one null leg plus one tachyon leg.  Two tachyon legs do not fail.  The statement r05.3(c) itself stands.

\subsection*{v4.9.0-r06 the hit chain at constant speed}
\label{subsec:v490-r06-chain}

\paragraph{Theoremlet v4.9.0-r06.2.}
Let hits occur along a line at speed \(v>1\), each emitting an identical light signal.  The energy of each signal comes
from the field at the hit, so nothing is carried faster than light.
\begin{enumerate}[leftmargin=2em]
\item[(a)] The signals superpose to a cone, the envelope of their Huygens wavelets, with \(\cos\vartheta=1/v\).  For
\(v=1/\bsym\) this gives \(\cos\vartheta=\bsym\), \(\vartheta=56.68^\circ\).
\item[(b)] The hits are spacelike separated.  Observers slower than \(1/v\) see them run in the direction of the
tachyon, observers faster than \(1/v\) see them run the other way, and the observer of speed exactly \(u=1/v\) sees all hits
at once.  For \(v=1/\bsym\) that observer moves at \(\bsym\), the matter-frame offset.
\item[(c)] The superluminal rapidity of r20.2 is \(\operatorname{artanh}v=\eta+i\pi/2\).  The boost to the observer of (b)
removes \(\eta\) and leaves \(i\pi/2\), the duality quarter turn of r07.2.  In that frame the tachyon is no longer a moving
object but a phase.  The plane wave emitted there perpendicular to the line appears in the original frame at
\(\cos\vartheta'=\bsym\), which is the cone of (a).
\item[(d)] With imaginary mass \(i\mu\), the tachyon at \(v=1/\bsym\) has energy \(\mu\sinh\eta\) and momentum
\(\mu\cosh\eta=1.197\,\mu\), the neighbour cross term of r18.
\end{enumerate}

\paragraph{Proof sketch.}
(a) Minimal arrival time over emission points; checked by a discrete Huygens sum.  (b)
\(\Delta t'=\gamma(\Delta t-u\Delta x)\) with \(\Delta t=\Delta x/v\).  (c) \(\tanh(x+i\pi/2)=\coth x\) and the aberration
formula with \(\vartheta=90^\circ\).  (d) \(E=\mu/\sqrt{v^2-1}\), \(p=\mu v/\sqrt{v^2-1}\).  \(\square\)

\paragraph{Reading.}
``The tachyon was already on its way before the experiment started'' is the frame dependence of (b).  There is an
observer for whom the whole path exists at one instant.  For \(v=1/\bsym\) this observer is the matter frame.  A
superluminal chain of locally triggered emissions is the situation of a light spot sweeping faster than light across a
screen.  It is permitted because each point emits by itself.

\subsection*{v4.9.0-r06 acceleration and infinite speed}
\label{subsec:v490-r06-acceleration}

\paragraph{Theoremlet v4.9.0-r06.3.}
\begin{enumerate}[leftmargin=2em]
\item[(a)] \emph{Constant proper acceleration.}  The tachyonic counterpart of hyperbolic motion is
\(t=R\cosh\sigma\), \(x=R\sinh\sigma\), with proper acceleration \(1/R\) everywhere.  Its speed \(v=\coth\sigma\) and its
coordinate acceleration \(dv/dt\) diverge at \(\sigma=0\), that is \(t=R\), \(x=0\).  The hit rate per unit time diverges
there as well, so the hits near that point arrive at once as a flash.  Beyond \(\sigma=0\) the parameter runs backward in
time.  By reinterpretation that part is a partner moving forward, and the point \((R,0)\) is a pair event.  A signal emitted
along the backward part is read as absorbed.
\item[(b)] \emph{Energy given to the hits.}  Write \(E=\mu\sinh\chi\), \(p=\mu\cosh\chi\), \(v=\coth\chi\).  If each hit
takes the energy \(\varepsilon\) from the tachyon, \(E\) falls and \(v\) rises.  After \(E_0/\varepsilon\) hits \(E=0\),
\(v=\infty\), \(p=\mu\).  Along the way the cone of r06.2 opens from \(56.68^\circ\) to \(90^\circ\), and the observer
seeing simultaneous hits slows from \(\bsym\) to rest.
\end{enumerate}
In both cases the infinite acceleration is a coordinate effect at a single point.  There the tachyon passes through infinite
speed, rapidity \(i\pi/2\) in the frame at hand.  The proper acceleration stays finite.

\paragraph{Proof sketch.}
(a) \(dx/dt=\coth\sigma\); \(dv/dt=-1/(R\sinh^3\sigma)\); the second derivative with respect to proper length has Minkowski
norm \(1/R\).  (b) \(p^2-E^2=\mu^2\) and \(v=p/E\).  Checked numerically, (b) with six hits from \(v=1/\bsym\).  \(\square\)

\paragraph{Reading.}
A tachyon that pays for the photons of its hits accelerates by itself and ends at infinite speed after finitely many
hits.  At that point the lab frame has become the frame of simultaneous hits, the cone has become a plane wave, and the
tachyon is the pure phase of r06.2(c).  Continuing past this point turns emission into absorption in the lab description.
This is the nearest the hit chain comes to the author's first picture of photons returning to the origin: a chain of
emissions meets a chain of absorptions at the point of infinite speed.

\subsection*{v4.9.0-r06 the Higgs field}
\label{subsec:v490-r06-higgs}

The main body assigns the Higgs sector the role of mediating between the compensated photon fix-sector and the mass
readout of localised branches (Definition~\ref{def:higgs-mediation}, Theorem~\ref{thm:higgs-role-refined}).

\paragraph{Theoremlet v4.9.0-r06.4.}
\begin{enumerate}[leftmargin=2em]
\item[(a)] In the electroweak theory the Higgs field gives the neutral pair \((W^3,B)\) the mass matrix
\(\tfrac{v_H^2}{4}\bigl(\begin{smallmatrix}g^2&-gg'\\-gg'&g'^2\end{smallmatrix}\bigr)\).  It has rank one: the photon is its massless
eigenvector \((\sin\theta_W,\cos\theta_W)\), the \(Z\) the massive one.  This is a port-weighting matrix in the sense of
r05.1, with weight \(0\) on the photon port and \(m_Z^2\) on the other.  For \(g=g'\) (\(\theta_W=45^\circ\)) it is exactly
\(0\cdot C_{\HC}+m_Z^2A\), and the photon is the compensator port.  The measured mixing angle differs from \(45^\circ\), so the
electroweak photon port is a compensator rotated by \(\theta_W-45^\circ\).
\item[(b)] At \(\phi=0\) the Higgs potential has negative mass squared.  The field there is tachyonic in the
field-theoretic sense, and mass arises when it condenses away from \(\phi=0\).
\end{enumerate}

\paragraph{Proof sketch.}
(a) The kernel of the matrix is spanned by \((g',g)\); for \(g=g'\) the matrix is \(\tfrac{v_H^2g^2}{2}A\).  (b) The
potential \(-\mu_H^2|\phi|^2+\lambda|\phi|^4\) has curvature \(-\mu_H^2\) at the origin.  \(\square\)

\paragraph{Reading.}
The Higgs part of the author's picture fits the main body in three respects.
\begin{itemize}[leftmargin=2em]
\item The photon is the direction the Higgs field does not weight, which is Theorem~\ref{thm:higgs-role-refined} in
electroweak form.
\item A field ``trying to create matter'' is a field leaving its tachyonic top.
\item A hit that yields a photon instead of a massive state lands in the unweighted port.
\end{itemize}
Three points need care.
\begin{itemize}[leftmargin=2em]
\item The doublet of r04 is \((E,B)\), the electroweak doublet is \((W^3,B)\).  The correspondence is formal.
\item A field with negative mass squared is an instability that grows; it is not a particle that carries signals faster
than light.  A superluminal chain of hits therefore needs triggers correlated in advance, as for the sweeping light spot,
and not a signal along the line.
\item The condensed Higgs vacuum is the minimum of its potential and supplies no energy.  The energy of the photons comes
either from the tachyon, case r06.3(b), or from the field while it is still rolling off its tachyonic top, that is during
symmetry breaking.
\end{itemize}
Comparing \(\theta_W\) or any other measured quantity with the values of this layer is a prediction and requires a
pre-registration note.

\subsection*{v4.9.0-r06 register entries}
\label{subsec:v490-r06-register}

\begin{center}
\small
\begin{tabular}{@{}>{\raggedright\arraybackslash}p{0.34\textwidth}>{\raggedright\arraybackslash}p{0.58\textwidth}@{}}
\toprule
statement & status after v4.9.0-r06\\
\midrule
r05.3(c), return photon & stands; one photon plus one tachyon leg leaves the gap \(d(1-1/v)\); two tachyon legs can close a loop (r06.1)\\
hit chain at speed \(1/\bsym\) & cone \(56.68^\circ\); all hits simultaneous for the observer at \(\bsym\), where the tachyon is the phase \(i\pi/2\) (r06.2)\\
accelerating tachyon & infinite speed at one point after finite time or finitely many hits; there a pair event, emission turns into absorption (r06.3)\\
Higgs mediation (Def.~\ref{def:higgs-mediation}) & electroweak mass matrix is a port weighting with massless photon port; compensator port for \(\theta_W=45^\circ\); formal correspondence (r06.4)\\
\bottomrule
\end{tabular}
\end{center}

\subsection*{v4.9.0-r06 ledger and guard update}
\label{subsec:v490-r06-ledger-guard}

\[
R_{\rm chain}^{v490\text{-}r06}=(0,0,1)
\]
for (loop condition decided, hit chain and acceleration worked out, physical status of the trigger and of the Higgs
correspondence open).

The layer spends no guard:
\[
R_{\rm glob}^{(v490\text{-}r06,+)}=(0,1,1,1,1),
\qquad
\Gamma_{\rm guard}^{v490\text{-}r06}=(1,1,1,1,1),
\qquad
\beta_4=1/256\ \text{unspent}.
\]
It does not imply exact finite \(C5\) zero closure, exact all-readout terminality,
\(C6/T2\), a global photon-clock/RH theorem, or global all-readout closure.

